\textbf{UC-ROZ-01: Zaksięgowanie wpłaty klienta}

\begin{longtable}{|p{4cm}|p{11cm}|}
\hline
\textbf{Element} & \textbf{Opis} \\
\hline

Identyfikator & UC-ROZ-01 \\
\hline

Nazwa & Zaksięgowanie wpłaty klienta \\
\hline

Opis &
Proces rejestracji wpływu środków oraz automatycznej klasyfikacji wpłaty jako zadatek, wpłata lub płatność końcowa. \\
\hline

Aktor główny & Księgowy / Bramka płatnicza \\
\hline

Aktorzy pomocniczy & Kontekst Sprzedaży i CRM, System Bankowy \\
\hline

Warunki wstępne &
W systemie istnieje aktywne zamówienie lub faktura proforma przypisana do klienta. \\
\hline

Warunki końcowe &
System rejestruje wpłatę, aktualizuje saldo rachunku oraz emituje zdarzenie dziedzinowe \textit{ZadatekZaksiegowany}. \\
\hline

Scenariusz główny &
\begin{enumerate}
\item System odbiera powiadomienie z bramki płatniczej lub Księgowy rejestruje wpłatę ręcznie w interfejsie użytkownika.
\item System identyfikuje zamówienie na podstawie identyfikatora płatności lub tytułu przelewu.
\item System weryfikuje kwotę wpłaty i przypisuje jej odpowiednią kategorię biznesową: Zadatek, Wpłata lub Płatność końcowa.
\item System rejestruje operację finansową w ewidencji księgowej.
\item System aktualizuje saldo przypisane do zamówienia.
\item System emituje zdarzenie dziedzinowe \textit{ZadatekZaksiegowany}, które odblokowuje dalszą realizację zamówienia w Module Sprzedaży.
\end{enumerate}
\\
\hline

Scenariusze alternatywne &
\begin{itemize}
\item [A1] Nierozpoznany tytuł przelewu: W kroku 2 system nie może powiązać wpłaty z istniejącym zamówieniem. System przenosi wpłatę do kolejki „Wpłaty do ręcznej weryfikacji”, a Księgowy dokonuje manualnego przypisania operacji.
\item [A2] Niedopłata: W kroku 3 system wykrywa, że wpłacona kwota jest niższa niż wymagane minimum zadatku. System księguje środki jako Wpłate, nie emituje zdarzenia \textit{ZadatekZaksiegowany} oraz wysyła powiadomienie do Handlowca.
\end{itemize}
\\
\hline

\end{longtable}